\chapter{Альтернативні сучасні теорії навчання на основі дофаміну}
\chaptermark{Інші теорії \nt{DOPA}}
\label{sec:dofaalt}
\begin{chaptergoals}
\item як \nt{DOPA}-нейрони з асиметричною відповіддю на помилки записують увесь розподіл винагороди;
\item як частина \nt{DOPA}-нейронів повідомляє важливість чи небезпеку, а не знак помилки;
\item як один провід несе швидку помилку, що вчить, і повільний рівень, що задає швидкість дій;
\item чому сигнал \nt{DOPA} --- джгут, а не одне число, і яка теорія ставить під сумнів саму помилку.
\end{chaptergoals}

\section{Що насправді йде проводом}

У розділі~\ref{sec:dofa} \nt{DOPA}-система була в нас одним проводом, яким іде одне число --- помилка передбачення винагороди $\delta$~\eqref{eq:ctd}. Ця картина пояснює багато, але що точніше вимірюють дофамін, то більше знаходять того, що в неї не вкладається. Одні нейрони відповідають на неприємне, як на винагороду; концентрація дофаміну росте, поки тварина біжить до мети, хоча нічого несподіваного не відбувається; різні \nt{DOPA}-нейрони поводяться по-різному.

Для інженера всі ці знахідки --- питання про одне: \emph{який формат має сигнал на широкомовній шині?} Одне число чи вектор? Один сигнал на проводі чи кілька, рознесених за частотою? Чи говорить він про майбутнє, як $\delta$, чи про минуле? Нижче --- шість сучасних відповідей; для кожної ми відокремимо виміряне від припущеного.

\section{Не одне число, а розподіл}

Помилка~\eqref{eq:ctd} вчить критика \emph{середнього} очікування винагороди. Але пів краплі соку напевно і крапля з імовірністю одна друга мають однакове середнє. Щоб їх розрізняти, потрібен увесь розподіл винагороди.

Візьмімо найпростішу задачу: після провісника надходить винагорода випадкового розміру $r$. Нехай \nt{DOPA}-нейронів багато, і кожен, $i$-й, обслуговує свою оцінку $V_i$, тож у момент винагороди його помилка дорівнює $\delta_i=r-V_i$. І нехай додатну помилку кожен враховує з вагою $\alpha_i^+$, від'ємну --- з вагою $\alpha_i^-$. Після кожної винагороди оцінка зсувається на
\begin{equation}
\Delta V_i \;=\; \eta\,\bigl(\alpha_i^+\,[\delta_i]_+ \;-\; \alpha_i^-\,[-\delta_i]_+\bigr).
\label{eq:asymmetric}
\end{equation}
Оцінка перестає змінюватися, коли в середньому додатні поправки врівноважують від'ємні:
\begin{empheq}[box=\keybox]{equation}
\alpha_i^+\,\bigl\langle[r-V_i]_+\bigr\rangle \;=\; \alpha_i^-\,\bigl\langle[V_i-r]_+\bigr\rangle ,
\qquad
\kappa_i \;=\; \frac{\alpha_i^+}{\alpha_i^++\alpha_i^-} .
\label{eq:expectile}
\end{empheq}
При $\kappa_i=1/2$ це звичайне середнє. При $\kappa_i>1/2$ нейрон --- «оптиміст»: його оцінка зсунута до верхнього краю розподілу, при $\kappa_i<1/2$ --- «песиміст».

Приклад: сік надходить з імовірністю $1/2$, розмір краплі $1$. Тоді~\eqref{eq:expectile} дає $\kappa_i\cdot\tfrac12(1-V_i)=(1-\kappa_i)\cdot\tfrac12V_i$, тобто $V_i=\kappa_i$ (рис.~\ref{fig:expectiles}). Набір нейронів із різними $\kappa_i$ записує розподіл винагороди так само, як набір квантилів записує розподіл у статистиці.

\begin{figure}[t]
\centering
\begin{tikzpicture}[>=Stealth,x=7cm,y=3cm]
\draw[->,gray!70!black] (-0.08,0) -- (1.15,0) node[right,black] {\small винагорода};
\draw[->,gray!70!black] (-0.08,0) -- (-0.08,0.75) node[left,black] {\small імовірність};
\fill[blue!25] (-0.03,0) rectangle (0.03,0.5);
\fill[blue!25] (0.97,0) rectangle (1.03,0.5);
\node[below,font=\small] at (0,0) {$0$};
\node[below,font=\small] at (1,0) {$1$};
\node[left,font=\scriptsize] at (-0.08,0.5) {$1/2$};
\foreach \k/\c/\lab/\h in {0.2/blue!60!black/{песиміст, $\kappa=0.2$}/0.62, 0.5/gray!50!black/{середнє, $\kappa=0.5$}/0.42, 0.8/red!70!black/{оптиміст, $\kappa=0.8$}/0.22}{
  \draw[very thick,\c] (\k,0) -- (\k,\h);
  \node[above,font=\scriptsize,\c] at (\k,\h) {\lab};}
\end{tikzpicture}
\caption{Розподіл винагороди записано набором \nt{DOPA}-нейронів. Винагорода $0$ або $1$ з імовірністю $1/2$ (сині стовпчики); нейрон із часткою $\kappa$ приходить до оцінки $V=\kappa$~\eqref{eq:expectile}.}
\label{fig:expectiles}
\end{figure}

Що виміряно: у різних \nt{DOPA}-нейронів різна «точка перевертання» --- розмір винагороди, за якого відповідь змінюється зі сплеску на провал, --- і нейрони з більшою асиметрією відповідей на додатні й від'ємні помилки перевертаються за більших винагород, як і вимагає~\eqref{eq:expectile}~\cite{Dabney2020}. З відповідей групи нейронів вдається відновити форму розподілу. Що це головна функція розкиду, а не побічний ефект, --- гіпотеза.

\begin{keyblock}
\begin{proposition}[Розподіл з асиметрії]
\label{prop:expectile}
Якщо \nt{DOPA}-нейрони різняться часткою $\kappa_i$, з якою вони враховують додатні помилки, їхні оцінки сходяться до різних точок~\eqref{eq:expectile} розподілу винагороди. Група нейронів передає не середнє, а розподіл, а отже --- ризик.
\end{proposition}
\end{keyblock}

\section{Не лише помилка, а й важливість}

За формулою помилки~\eqref{eq:ctd} неприємна подія --- удар, гіркий смак --- має спричиняти провал: вийшло гірше, ніж очікувалося. Частина \nt{DOPA}-нейронів так і робить, а інша відповідає на неприємне \emph{сплеском}, як на винагороду~\cite{Matsumoto2009}. Ці нейрони повідомляють не знак помилки, а те, що сталося щось \emph{важливе}: несподіване, сильне, що потребує уваги~\cite{BrombergMartin2010}.

Інженерові зручно думати про це як про два сигнали. Перший --- знакова помилка $\delta$: куди змінювати ваги. Другий --- її модуль $|\delta|$ або новизна події: \emph{наскільки сильно} їх змінювати, тобто швидкість навчання; після несподіваної події навчатися треба швидше. За змістом він близький до норадреналінового коефіцієнта підсилення з розділу~\ref{sec:neurontypes}.

Є й третій варіант: окремий провід для окремої осі. \nt{DOPA}-нейрони, що йдуть до однієї з областей вибору дій, відповідають не на винагороду, а на новизну і силу подразника і потрібні, щоб тварина навчилася уникати небезпечного~\cite{Menegas2018}: проводом іде не «краще чи гірше», а «небезпечно чи ні».

Що виміряно: різні групи \nt{DOPA}-нейронів по-різному відповідають на неприємне і на нове, і різні виходи вчать різного. Як саме мозок користується сигналом важливості --- як швидкістю навчання, як сигналом уваги чи як окремою помилкою за віссю небезпеки, --- обговорюється.

\section{Не лише вчити, а й підганяти}

Дофамін має й негайну дію: що його більше, то охочіше і швидше тварина діє. Як поєднати це з навчанням?

Інженерна відповідь --- \emph{розділення за частотою}. Швидкі сплески і провали тривалістю в сотні мілісекунд несуть $\delta$ і вчать. Повільний рівень, що змінюється за хвилини, несе іншу величину --- середню винагороду за одиницю часу $\bar R$~\cite{Niv2007}. Нехай дія, виконана за час $\tau_a$, потребує зусилля $C/\tau_a$: що швидше, то дорожче. Натомість кожна секунда зволікання коштує $\bar R$ --- стільки винагороди можна було б за неї отримати. Сумарна ціна $C/\tau_a+\bar R\,\tau_a$ мінімальна при
\begin{empheq}[box=\keybox]{equation}
\tau_a^{*} \;=\; \sqrt{\frac{C}{\bar R}} .
\label{eq:vigor}
\end{empheq}
Що багатша обстановка, то дорожчий час і то швидше вигідно діяти: якщо $\bar R$ подвоїлася, час дії зменшується в $\sqrt2\approx1.4$ раза. Зауважте, що $\bar R$ --- та сама очікувана винагорода, яку віднімали в розділі~\ref{sec:dofa}.

Викид дофаміну в місці призначення може змінюватися і без зміни частоти імпульсів: його підлаштовують місцеві сигнали біля самих виходів~\cite{Mohebi2019}. Але повільні складові є і в самих імпульсах: у дослідах наступного підрозділу плавне зростання під час наближення до винагороди видно і в частоті \nt{DOPA}-нейронів~\cite{Kim2020}. Отже, повільний рівень складається з двох джерел --- частоти імпульсів і місцевого регулювання викиду, --- і яке з них головне, схоже, залежить від виходу і від задачі.

\begin{keyblock}
\begin{proposition}[Два сигнали на одному проводі]
\label{prop:multiplex}
\nt{DOPA}-система передає щонайменше дві величини, рознесені в часі: швидку помилку $\delta$, яка вчить, і повільний рівень, який задає ціну часу~\eqref{eq:vigor} і тим самим швидкість дій. Виміряно, що швидка і повільна складові поводяться по-різному; що повільна дорівнює саме $\bar R$, --- гіпотеза.
\end{proposition}
\end{keyblock}

\section{Зростання під час наближення}

Якщо тварина біжить коридором до далекої винагороди, концентрація дофаміну в області вибору дій плавно росте протягом секунд, поки мета наближається~\cite{Hamid2016}. За розділом~\ref{sec:dofa} цього не мало б бути: усе йде за планом, і $\delta$ має дорівнювати нулю.

Перше пояснення: це зростання --- не $\delta$, а саме очікування $V$, сигнал «мета близько, варто постаратися» з попереднього підрозділу~\cite{Hamid2016}. Друге пояснення зберігає $\delta$~\cite{Kim2020}. Якщо очікування правильне, то за горизонту $\tau_\gamma$ на шляху до винагороди $R$ воно росте як $V=Re^{-(T-t)/\tau_\gamma}$, і за формулою помилки~\eqref{eq:ctd} $\dd V/\dd t-V/\tau_\gamma=0$. Але нехай здалеку очікування занижене, а в міру наближення оцінка уточнюється. Тоді очікування росте крутіше, скажімо як $V=Re^{-(T-t)/\tau_v}$ з $\tau_v<\tau_\gamma$, і
\begin{empheq}[box=\keybox]{equation}
\delta(t) \;=\; \frac{\dd V}{\dd t}-\frac{V}{\tau_\gamma} \;=\; V\Bigl(\frac{1}{\tau_v}-\frac{1}{\tau_\gamma}\Bigr) \;>\; 0 .
\label{eq:ramp}
\end{empheq}
Помилка додатна і росте разом із $V$ --- те саме плавне зростання (рис.~\ref{fig:ramp}). При $\tau_\gamma=4\,\text{с}$ і $\tau_v=1\,\text{с}$ виходить $\delta=0.75\,V$ на секунду. Цю версію перевіряли у віртуальному коридорі, миттєво переносячи тварину ближче до мети: дофамін відповідав сплеском, як і належить похідній, а не плавному рівню~\cite{Kim2020}.

\begin{figure}[t]
\centering
\begin{tikzpicture}[>=Stealth,x=1.6cm,y=2.6cm]
\draw[->,gray!70!black] (-4.2,0) -- (0.5,0) node[below left,black] {\small $t$};
\draw[->,gray!70!black] (-4.2,0) -- (-4.2,1.2);
\foreach \x/\l in {-4/{$T-4$},-2/{$T-2$},0/{$T$}}{\draw[gray!60!black] (\x,-0.03) -- (\x,0.03); \node[below,font=\scriptsize] at (\x,0) {\l\,с};}
\draw[thick,densely dashed,gray!60!black] plot[domain=-4:0,samples=40] (\x,{exp(\x/4)});
\draw[very thick,blue!60!black] plot[domain=-4:0,samples=60] (\x,{exp(\x)});
\draw[very thick,red!70!black] plot[domain=-4:0,samples=60] (\x,{0.75*exp(\x)});
\node[anchor=south east,font=\small,gray!50!black] at (-1.3,0.77) {$V$ при $\tau_v=\tau_\gamma$: $\delta=0$};
\node[right,font=\small,blue!60!black] at (0.05,1.0) {$V$ при $\tau_v=1\,\text{с}$};
\node[right,font=\small,red!70!black] at (0.05,0.72) {$\delta$ за~\eqref{eq:ramp}};
\end{tikzpicture}
\caption{Зростання дофаміну під час наближення до винагороди як помилка передбачення, $\tau_\gamma=4\,\text{с}$. Штрихова лінія --- очікування, що росте рівно так, як вимагає горизонт ($\delta=0$); синя --- очікування, що росте крутіше; червона --- помилка~\eqref{eq:ramp}.}
\label{fig:ramp}
\end{figure}

Що виміряно: плавне зростання є --- і в концентрації дофаміну, і в частоті імпульсів \nt{DOPA}-нейронів~\cite{Kim2020}, --- і миттєві стрибки обстановки спричиняють стрибки дофаміну. Яке з двох пояснень правильне --- чи обидва правильні на різних виходах, --- обговорюється.

\section{Погляд назад}

Усі теорії вище дивляться \emph{вперед}. Одна із сучасних теорій перевертає напрямок~\cite{Jeong2022}. Нехай мозок навчається не передбачати винагороду за провісником, а, отримавши винагороду, \emph{озиратися}: яка подія частіше за інші їй передувала? Міра такого зв'язку --- різниця двох імовірностей:
\begin{empheq}[box=\keybox]{equation}
M \;=\; P(\text{провісник перед винагородою}) \;-\; P(\text{провісник у випадковому вікні}) .
\label{eq:retro}
\end{empheq}
Якщо провісник часто буває і без винагороди, другий доданок великий, і зв'язок слабкий. У цій теорії дофамін повідомляє не помилку передбачення, а те, наскільки дана подія --- імовірна \emph{причина} винагороди.

Механізм озирання потребує того самого, що й правило трьох множників розділу~\ref{sec:dofa}: кожна подія має лишати слід, що тане, щоб у момент винагороди можна було прочитати, що було перед нею. Різниця не в синапсі, а в тому, що передає \nt{DOPA}.

У дослідах, де теорія~\eqref{eq:retro} і помилка~\eqref{eq:ctd} передбачають різне, викид дофаміну в низці випадків ішов за першою~\cite{Jeong2022}. Але в інших дослідах відповідь дофаміну під час навчання поступово зсувалася від моменту винагороди до моменту провісника --- як вимагає навчання за помилкою~\eqref{eq:ctd}~\cite{Amo2022}. Яка теорія описує мозок, поки невідомо.

\section{Помилка не лише у винагороді}

Останнє розширення --- про те, \emph{про що} помилка. Якщо замінити очікуваний сік іншим, так само цінним, але іншого смаку, цінність не змінилася, і помилка~\eqref{eq:ctd} дорівнює нулю. Проте \nt{DOPA}-нейрони на таку заміну відповідають~\cite{Takahashi2017}. А штучно викликані сплески дофаміну можуть навчити тварину зв'язку між двома нейтральними подразниками, взагалі без винагороди~\cite{Sharpe2017}. Схоже, \nt{DOPA} повідомляє про помилку передбачення не лише винагороди, а й того, \emph{що} станеться.

Формально це просте узагальнення~\eqref{eq:ctd}: замість одного потоку винагороди $r$ береться набір ознак того, що відбувається, $\varphi_k$ --- смак, місце, звук, --- і для кожної своє очікування $M_k$ і своя помилка~\cite{Gardner2018}:
\begin{empheq}[box=\keybox]{equation}
\delta_k(t) \;=\; \varphi_k(t) \;+\; \frac{\dd M_k}{\dd t} \;-\; \frac{M_k}{\tau_\gamma} ,
\qquad k=1,\dots,K .
\label{eq:vectortd}
\end{empheq}
Винагорода --- лише одна з ознак, і вектор помилок учить не тільки того, що добре, а й того, що за чим іде, --- моделі світу.

З цим узгоджується неоднорідність \nt{DOPA}-нейронів: в одній задачі різні нейрони несуть різні величини --- одні пов'язані з винагородою, інші з рухом, треті з тим, куди спрямована увага~\cite{Engelhard2019}.

\begin{keyblock}
\begin{proposition}[Джгут замість проводу]
\label{prop:bundle}
Сигнал \nt{DOPA}-системи --- не одне число. Він розкладений за нейронами (розподіл~\eqref{eq:expectile}, різні ознаки~\eqref{eq:vectortd}, окрема вісь небезпеки) і за часом (швидка помилка і повільний рівень~\eqref{eq:vigor}). Скалярна помилка~\eqref{eq:ctd} --- перше наближення до цього сигналу, як суматор із порогом --- перше наближення до нейрона.
\end{proposition}
\end{keyblock}

\section{Підсумок}

\begin{table}[t]
\centering
\footnotesize
\setlength{\tabcolsep}{4pt}
\renewcommand{\arraystretch}{1.15}
\begin{tabular}{>{\raggedright}p{2.5cm}>{\raggedright}p{3.2cm}>{\raggedright}p{3.6cm}>{\raggedright\arraybackslash}p{4.2cm}}
\toprule
Теорія & Що йде проводом & Головне вимірювання & Що відомо \\
\midrule
помилка передбачення (розд.~\ref{sec:dofa}) & скаляр $\delta$~\eqref{eq:ctd} & сплеск, перенесення на провісник, провал & виміряно; перше наближення \\
розподіл & набір $\delta_i$ з різною асиметрією & різні точки перевертання в різних нейронів & узгодженість із дослідом є; роль розкиду --- гіпотеза \\
важливість & $|\delta|$, новизна; вісь небезпеки & сплески на неприємне в частини нейронів & виміряно; як використовується --- обговорюється \\
ціна часу & повільний рівень $\bar R$ & дофамін пришвидшує дії; повільна складова --- і у викиді біля виходів, і в частоті імпульсів & розділення швидкого і повільного виміряно; $\bar R$ --- гіпотеза \\
зростання під час наближення & $V$ або $\delta$ при уточненні оцінки~\eqref{eq:ramp} & плавне зростання, стрибки при перенесенні в коридорі & зростання виміряно; пояснення обговорюється \\
погляд назад & міра причинного зв'язку~\eqref{eq:retro} & досліди, де передбачення двох теорій розходяться & результати суперечать одне одному \\
вектор помилок & $\delta_k$ за ознаками~\eqref{eq:vectortd} & відповідь на зміну смаку; навчання без винагороди & виміряно; векторна форма --- гіпотеза \\
\bottomrule
\end{tabular}
\caption{Що передає \nt{DOPA}-система: стандартна картина розділу~\ref{sec:dofa} та її сучасні розширення.}
\label{tab:dofaalt}
\end{table}

Теорії таблиці~\ref{tab:dofaalt} не виключають одна одну: майже всі зберігають помилку передбачення як основу і розширюють її формат. Усерйоз конкурує з нею лише погляд назад, і цю суперечку розв'яжуть досліди. Для інженера висновок простий: ціна широкомовлення з твердження~\ref{prop:broadcastcost} падає, якщо провід насправді --- багато проводів із різними адресатами.

\section{Задачі}

\begin{exercise}[\lvlA]
\label{ex:07:1}
\emph{Точки розподілу для однієї краплі.} Сік розміру $1$ надходить з імовірністю $p=0.2$, інакше винагорода $0$. Знайдіть за~\eqref{eq:expectile} $V_i$ для $\kappa_i=0.2$, $0.5$, $0.8$. Чому дорівнює медіана цієї винагороди, і чим точка~\eqref{eq:expectile} відрізняється від квантиля?
\end{exercise}

\begin{exercise}[\lvlB]
\label{ex:07:2}
\emph{Розподіл за двома нейронами.} Винагорода дорівнює $0$ або $x$, імовірність $x$ дорівнює $p$. Два нейрони з правилом~\eqref{eq:asymmetric} повідомили $V(1/2)=0.25$ і $V(0.8)=0.5$. Відновіть $p$, $x$ і стандартне відхилення винагороди. Що повідомить нейрон із $\kappa=0.2$?
\end{exercise}

\begin{exercise}[\lvlB]
\label{ex:07:3}
\emph{Моделювання: розподіл з асиметрії.} Змоделюйте дев'ять \nt{DOPA}-нейронів з $\kappa_i=0.1,\dots,0.9$ і правилом~\eqref{eq:asymmetric}, $\alpha_i^+=2\kappa_i$, $\alpha_i^-=2(1-\kappa_i)$, за винагороди, рівномірної на $[0,1]$. Як розкид $V_i$ залежить від $\eta$, і яке $\eta$ потрібне, щоб відрізнити цей розподіл від двох рівноймовірних крапель $0.25$ і $0.75$?
\end{exercise}

\begin{exercise}[\lvlA]
\label{ex:07:4}
\emph{Ціна часу.} (а)~Виведіть~\eqref{eq:vigor} і знайдіть сумарну ціну в оптимумі. (б)~Мозок оцінив $\bar R$ з помилкою в $k$ разів. Знайдіть відносну переплату при $k=2$ і $k=4$. Наскільки точно повільний рівень дофаміну має повідомляти $\bar R$?
\end{exercise}

\begin{exercise}[\lvlB]
\label{ex:07:5}
\emph{Сплеск при перенесенні.} Очікування росте за~\eqref{eq:ramp} з $\tau_v=1\,\text{с}$, $\tau_\gamma=4\,\text{с}$; тварину миттєво переносять із точки $T-2\,\text{с}$ у $T-1\,\text{с}$. Знайдіть площу сплеску $\delta$ у частках $R$ і скільки секунд рампи до перенесення він замінює. Що покаже перенесення, якщо дофамін повідомляє саме $V$?
\end{exercise}

\begin{exercise}[\lvlB]
\label{ex:07:6}
\emph{Проєктування: два сигнали на проводі.} Проводом \nt{DOPA} ідуть рівень, що росте до $s=1/60$ імпульсу на секунду за секунду, і події по $A=3$ імпульси. Синапс зі слідом $\tau_e=1\,\text{с}$ віднімає від частоти її копію, згладжену з $\tau_f$. За яких $\tau_f$ втрачається не більше $10\,\%$ події, а внесок рівня за час сліду менший за $10\,\%$ внеску події?
\end{exercise}

\begin{exercise}[\lvlB]
\label{ex:07:7}
\emph{Проєктування досліду: погляд назад проти помилки.} У вікні провісник буває з імовірністю $0.1$, після нього винагорода --- з імовірністю $0.8$. Порівняйте $\Delta P=P(\text{вин.}\mid\text{пров.})-P(\text{вин.}\mid\text{нема})$ і $M$ з~\eqref{eq:retro}, якщо (а)~додати винагороди без провісника, $0.08$ на вікно; (б)~подвоїти частоту провісника без винагород.
\end{exercise}

\begin{exercise}[\lvlC]
\label{ex:07:8}
\emph{Джгут і ціна широкомовлення.} У моделі задачі~\ref{ex:06:8} $N$ нейронів поділено на $K$ груп із власними проводами, і в кожен провід просочується частка $\rho$ чужих помилок. Покажіть, що стала навчання дорівнює $N/K+2+\rho^2N(1-1/K)$. Скільки спроб вона дає при $K\to\infty$, $N=10^{10}$ і $\rho=0.01$?
\end{exercise}
